% ============================================================
%  Introducing Python, chapter 1: Introduction
% ============================================================
\chapter{Introduction}

This chapter shows what a Python program looks like, how to install Python, and
how to run code in two ways: typed line by line, or saved in a file. It ends with
a realistic script that reads JSON, the data format almost every AI API speaks.
You do not need to understand every line yet. You need to see the shape of the thing.

% ------------------------------------------------------------
\section{Mysteries}

A knitting pattern and a cooking recipe both look like nonsense to an outsider.
Yet both are \emph{little languages}: short vocabularies, strict rules, steps in order.
A program is the same kind of thing, written for a computer. The table maps the parts
of a recipe to the parts of a Python program.

\begin{center}
\begin{tabularx}{\linewidth}{@{}l X X@{}}
\toprule
\textbf{Feature} & \textbf{In a recipe} & \textbf{In Python} \\
\midrule
Vocabulary & ``fold'', ``sift'', abbreviations & keywords and names: \code{for}, \code{print} \\
Syntax & what may be written, and where & the grammar rules of the language \\
Sequence & steps run top to bottom & statements run top to bottom \\
Loop & ``repeat for each ball'' & \code{for} and \code{while} \\
Call another recipe & ``see the dough recipe'' & a \textbf{function} call \\
Assumed context & you know what an oven is & the standard library, the OS \\
Data & flour, potatoes & values: numbers, text, lists \\
Tools & bowl, pan & the interpreter, an editor, a terminal \\
Result & dinner & output, a file, a changed value \\
\bottomrule
\end{tabularx}
\end{center}

\begin{keyidea}[: programming is learnable]
Programming is mostly learning a small vocabulary and a few rules. Python keeps both
small, which is one reason it dominates data science and AI work.
\end{keyidea}

% ------------------------------------------------------------
\section{Little Python Programs}

Python is a general-purpose language created by Guido van Rossum in 1991. It is
designed to be easy to read, it runs tiny scripts and huge systems, and it is the main
language of modern machine learning. Three small programs show most of its flavor.

\subsection{A loop}

\begin{pycode}
# file: loop1.py
for step in 3, 2, 1, "go":
    print(step)
\end{pycode}
\begin{rawcode}
3
2
1
go
\end{rawcode}

\code{for} takes each value in turn, puts it in the name \code{step}, and runs the
indented line. Notice that the values are mixed: numbers and text. Python does not mind.

\subsection{A list: values by position}

A \textbf{list} keeps values in order inside square brackets. You reach one by its
\textbf{offset}: the distance from the start. The first item is at offset 0.

\begin{pycode}
models = ["haiku", "sonnet", "opus"]
print(models[1])
\end{pycode}

%MD
%```text
%  offset:     0         1          2
%           +---------+----------+--------+
%  models = | "haiku" | "sonnet" | "opus" |
%           +---------+----------+--------+
%  offset:    -3        -2         -1
%```
%ENDMD
\begin{center}
\begin{tikzpicture}[font=\ttfamily\small]
  \foreach \k/\v in {0/haiku,1/sonnet,2/opus} {
    \node[draw=cblue,fill=cblue!8,minimum width=2.3cm,minimum height=8mm] at (2.6*\k,0) {"\v"};
    \node[font=\sffamily\scriptsize,text=cgreen] at (2.6*\k,0.75) {offset \k};
  }
  \node[font=\sffamily\scriptsize,text=cred] at (0,-0.75) {offset -3};
  \node[font=\sffamily\scriptsize,text=cred] at (2.6,-0.75) {offset -2};
  \node[font=\sffamily\scriptsize,text=cred] at (5.2,-0.75) {offset -1};
  \node[anchor=east] at (-1.3,0) {models =};
\end{tikzpicture}
\end{center}

The output of the program above is \code{sonnet}: offset 1 is the \emph{second} item.
Think in offsets, not positions.

\subsection{A dictionary: values by name}

A \textbf{dictionary} stores \emph{key}: \emph{value} pairs in curly braces. You look a
value up by its key instead of by a number.

\begin{pycode}
context_limit = {"small": 8_000, "medium": 32_000, "large": 128_000}
size = "medium"
print(size, "allows", context_limit[size], "tokens")
\end{pycode}

(The numbers are invented for the demo.) Output: \code{medium allows 32000 tokens}.

%MD
%```mermaid
%flowchart LR
%    K1["'small'"] --> V1["8000"]
%    K2["'medium'"] --> V2["32000"]
%    K3["'large'"] --> V3["128000"]
%```
%ENDMD
\begin{center}
\begin{tikzpicture}[font=\ttfamily\small,k/.style={draw=cgold,fill=cgold!10,minimum width=2cm,minimum height=7mm},
  v/.style={draw=cgreen,fill=cgreen!10,minimum width=2cm,minimum height=7mm}]
  \foreach \i/\a/\b in {0/small/8000,1/medium/32000,2/large/128000} {
    \node[k] (k\i) at (0,-0.9*\i) {"\a"};
    \node[v] (v\i) at (4,-0.9*\i) {\b};
    \draw[->,thick,cgray] (k\i) -- (v\i);
  }
  \node[font=\sffamily\scriptsize,text=cgray] at (0,0.7) {keys};
  \node[font=\sffamily\scriptsize,text=cgray] at (4,0.7) {values};
\end{tikzpicture}
\end{center}

\begin{typecheck}[: square brackets do three different jobs]
\begin{center}
\begin{tabular}{@{}l l l@{}}
\toprule
\textbf{You write} & \textbf{What it is} & \textbf{You get} \\
\midrule
\code{["a", "b"]} & list literal & a \code{list} \\
\code{models[1]} & subscript with an int offset & one item from the list \\
\code{context\_limit["small"]} & subscript with a key & one value from the dict \\
\bottomrule
\end{tabular}
\end{center}
\medskip
Using the wrong thing fails loudly:
\begin{pycode}
>>> models = ["haiku", "sonnet", "opus"]
>>> models[3]
Traceback (most recent call last):
  ...
IndexError: list index out of range
>>> {"small": 8_000}["tiny"]
Traceback (most recent call last):
  ...
KeyError: 'tiny'
\end{pycode}
\end{typecheck}

% ------------------------------------------------------------
\section{Setup}

To write Python you need two tools: a \textbf{terminal} (a window where you type
commands) and a \textbf{text editor} (plain text, not a word processor). Windows,
macOS and Linux all have both. IDEs such as VS Code bundle them; they come in a
later chapter.

\begin{trap}
Never write code in Word or similar. They swap straight quotes for curly ones
and add formatting. Python then fails with a confusing \code{SyntaxError}. Use a code editor.
\end{trap}

\subsection{Install Python}

Download the latest release from \code{python.org} (Downloads, then your platform)
and follow the installer. On Windows, tick the option that adds Python to \code{PATH}.
The computer may already have an old Python; install a current one anyway.

\subsection{Upgrade Python}

Versions are numbered \textbf{x.y.z}: a major number, a minor number, a bug-fix number.
Python 3 has been the major version for years and no Python 4 is planned. A new
minor version (3.12, 3.13, \ldots) arrives about once a year. Check yours:

\begin{rawcode}
$ python -V
Python 3.12.10
\end{rawcode}

On Windows the launcher may be \code{py -V}. This series uses Python 3.12. Newer minor
versions add features; code written for 3.12 normally runs on later ones.

\subsection{Run Python programs}

%MD
%```mermaid
%flowchart LR
%    T[terminal] -->|"python"| R["REPL: type a line, see the result"]
%    T -->|"python hello.py"| F["run the whole file, then exit"]
%    E[editor] -->|save| H[hello.py]
%    H --> F
%```
%ENDMD
\begin{center}
\begin{tikzpicture}[font=\sffamily\small,n/.style={draw=cblue,fill=cblue!8,rounded corners=2pt,inner sep=5pt},
  >=Stealth]
  \node[n] (t) at (0,0) {terminal};
  \node[n,fill=cgreen!10,draw=cgreen,align=center] (r) at (5.2,0.9) {interactive interpreter (REPL)\\[-1pt]{\scriptsize one line at a time}};
  \node[n,fill=cgold!10,draw=cgold,align=center] (f) at (5.2,-0.9) {run a file\\[-1pt]{\scriptsize whole program, then exit}};
  \node[n,draw=cgray,fill=cgray!8] (h) at (0,-1.8) {editor: hello.py};
  \draw[->,thick] (t) -- node[above,sloped,font=\ttfamily\scriptsize]{python} (r);
  \draw[->,thick] (t) -- node[below,sloped,font=\ttfamily\scriptsize]{python hello.py} (f);
  \draw[->,thick,cgray] (h) -- (f);
\end{tikzpicture}
\end{center}

\textbf{Interactive interpreter.} Type \code{python} alone. You get a prompt
\code{>{}>{}>}. Type one line of Python, press Enter, see the result at once. This is the
fastest way to test an idea.

\textbf{A file.} Save code in a file ending in \code{.py}, then run
\code{python hello.py}. Use files for anything longer than a few lines.

One difference surprises beginners. The interpreter \emph{echoes} the value of an
expression. A file does not. In a file you must call \code{print()}.

\begin{pycode}
>>> 6 * 7
42
>>> print(6 * 7)
42
>>> "hi"
'hi'
>>> print("hi")
hi
\end{pycode}

Note the quotes: the echo shows the value with quotes (a \emph{representation}),
while \code{print()} shows the text for humans.

\begin{typecheck}[: \texttt{print()} returns nothing]
\code{print()} displays text and returns \code{None}. If you write
\code{x = print("hi")}, then \code{x} is \code{None}, not the text. Use \code{print}
to show, not to produce a value.
\end{typecheck}

% ------------------------------------------------------------
\section{Built-In Python Features}

Where does code come from? Three layers, each wider than the last.

%MD
%```mermaid
%flowchart TB
%    A["Built-ins: print, len, int, str ... always there"] --> B["Standard library: json, math, pathlib ... installed with Python, use import"]
%    B --> C["Third-party packages: anthropic, openai, requests ... pip install, then import"]
%```
%ENDMD
\begin{center}
\begin{tikzpicture}[font=\sffamily\small]
  \draw[cgray,fill=cpurple!8,rounded corners=3pt] (0,0) rectangle (11,3.1);
  \draw[cgray,fill=cteal!10,rounded corners=3pt] (0.3,0.3) rectangle (7.4,2.3);
  \draw[cgray,fill=cgold!15,rounded corners=3pt] (0.6,0.6) rectangle (4.0,1.5);
  \node[align=center] at (2.3,1.05) {\textbf{built-ins}\\[-1pt]{\scriptsize\ttfamily print, len, int}};
  \node[align=center] at (5.7,1.4) {\textbf{standard library}\\[-1pt]{\scriptsize\ttfamily json, math, pathlib}\\[-1pt]{\scriptsize import it}};
  \node[align=center] at (9.2,1.55) {\textbf{third party}\\[-1pt]{\scriptsize\ttfamily anthropic, openai}\\[-1pt]{\scriptsize pip install, then import}};
\end{tikzpicture}
\end{center}

\begin{center}
\begin{tabularx}{\linewidth}{@{}l X X@{}}
\toprule
\textbf{Layer} & \textbf{How you get it} & \textbf{Example} \\
\midrule
Built-ins & nothing to do & \code{len("abc")} \\
Standard library & \code{import name} & \code{import math} \\
Third-party packages & \code{pip install name}, then import & \code{import anthropic} \\
\bottomrule
\end{tabularx}
\end{center}

\begin{pycode}
>>> import math
>>> math.sqrt(16)
4.0
>>> from json import dumps
>>> dumps({"ok": True})
'{"ok": true}'
\end{pycode}

Installing third-party packages is done in a terminal, ideally inside a virtual
environment (a private folder of packages per project):

\begin{rawcode}
$ pip install anthropic
\end{rawcode}

(\code{anthropic} is a third-party package; it is not needed to read this chapter.)

\begin{gamerule}
One virtual environment per project. Never install project packages into the
system Python. It keeps versions from clashing.
\end{gamerule}

% ------------------------------------------------------------
\section{A Bigger Example}

Here is a short script that does real work: it takes the JSON text an AI API might
send back, turns it into Python data, and pulls out two facts. It uses an import, a
multi-line string, a function call, nested lookups, an f-string and error handling.
(The JSON is canned, so the script runs offline.)

\begin{pycode}
# file: reply.py
import json

raw = """{
  "model": "demo-1",
  "usage": {"input": 12, "output": 40},
  "reply": {"text": "Hi there"}
}"""

data = json.loads(raw)
try:
    text = data["reply"]["text"]
    total = data["usage"]["input"] + data["usage"]["output"]
    print(f"{data['model']} said {text!r} using {total} tokens")
except KeyError as err:
    print("missing key:", err)
\end{pycode}
\begin{rawcode}
demo-1 said 'Hi there' using 52 tokens
\end{rawcode}

What each part does:

\begin{enumerate}[itemsep=1pt]
\item \code{import json}: make the standard-library JSON module available.
\item The triple-quoted string \code{raw} holds JSON \emph{text}. It is a \code{str}.
\item \code{json.loads(raw)} parses text into Python data. Here you get a \code{dict}.
\item \code{try:} run the indented block. If a key is missing, jump to \code{except}.
\item \code{data["reply"]["text"]} follows two layers of dictionaries.
\item The f-string \code{f"..."} inserts values inside curly braces.
\item \code{except KeyError as err} runs only if a key was absent; \code{err} holds the key.
\end{enumerate}

%MD
%```mermaid
%flowchart TD
%    D["data (dict)"] --> M["'model' -> 'demo-1'"]
%    D --> U["'usage' (dict)"]
%    D --> R["'reply' (dict)"]
%    U --> UI["'input' -> 12"]
%    U --> UO["'output' -> 40"]
%    R --> RT["'text' -> 'Hi there'"]
%```
%ENDMD
\begin{center}
\begin{tikzpicture}[font=\ttfamily\scriptsize,
  d/.style={draw=cblue,fill=cblue!8,rounded corners=2pt,inner sep=3pt},
  l/.style={draw=cgreen,fill=cgreen!10,rounded corners=2pt,inner sep=3pt},
  >=Stealth]
  \node[d] (r) at (5,0) {data (dict)};
  \node[l] (m) at (0.3,-1.4) {"model": "demo-1"};
  \node[d] (u) at (5,-1.4) {"usage" (dict)};
  \node[d] (p) at (9.6,-1.4) {"reply" (dict)};
  \node[l] (i) at (3.4,-2.8) {"input": 12};
  \node[l] (o) at (6.6,-2.8) {"output": 40};
  \node[l] (t) at (9.6,-2.8) {"text": "Hi there"};
  \foreach \a/\b in {r/m,r/u,r/p,u/i,u/o,p/t} \draw[->,thick,cgray] (\a) -- (\b);
\end{tikzpicture}
\end{center}

If the key is missing, the script reports it instead of crashing:

\begin{pycode}
>>> import json
>>> data = json.loads('{"model": "demo-1"}')
>>> try:
...     data["reply"]["text"]
... except KeyError as err:
...     print("missing key:", err)
...
missing key: 'reply'
\end{pycode}

%MD
%```mermaid
%flowchart TD
%    S[start try block] --> Q{"every key present?"}
%    Q -->|yes| OK["run all lines, skip except"]
%    Q -->|"no: KeyError"| EX["jump to except block"]
%    OK --> END[continue after]
%    EX --> END
%```
%ENDMD
\begin{center}
\begin{tikzpicture}[font=\sffamily\small,n/.style={draw=cblue,fill=cblue!8,rounded corners=2pt,inner sep=4pt},>=Stealth]
  \node[n] (s) at (0,0) {enter \code{try}};
  \node[n,rounded corners=9pt,draw=cgold,fill=cgold!10] (q) at (4,0) {key present?};
  \node[n,draw=cgreen,fill=cgreen!10] (ok) at (8,0.8) {run all lines};
  \node[n,draw=cred,fill=cred!8] (ex) at (8,-0.8) {jump to \code{except}};
  \node[n,draw=cgray,fill=cgray!8] (e) at (11.2,0) {carry on};
  \draw[->,thick] (s) -- (q);
  \draw[->,thick,cgreen] (q) -- node[above,sloped,font=\scriptsize]{yes} (ok);
  \draw[->,thick,cred] (q) -- node[below,sloped,font=\scriptsize]{no} (ex);
  \draw[->,thick] (ok) -- (e);
  \draw[->,thick] (ex) -- (e);
\end{tikzpicture}
\end{center}

\begin{keyidea}[: indentation is syntax]
Many languages mark blocks with braces. Python uses \textbf{indentation}. The lines
indented under \code{try:} belong to it; the unindented line ends it. Use four spaces.
Mixing tabs and spaces causes errors.
\end{keyidea}

\begin{trap}
A bare \code{except:} (with no error type) hides every bug, even a typo.
Name the error you expect: \code{except KeyError}.
\end{trap}

\begin{typecheck}[: what \texttt{json.loads} gives you]
\begin{center}
\begin{tabular}{@{}l l@{}}
\toprule
\textbf{JSON} & \textbf{Python} \\
\midrule
object \code{\{...\}} & \code{dict} \\
array \code{[...]} & \code{list} \\
string & \code{str} \\
number & \code{int} or \code{float} \\
\code{true} / \code{false} & \code{True} / \code{False} \\
\code{null} & \code{None} \\
\bottomrule
\end{tabular}
\end{center}
\medskip
A dict uses \emph{subscripts}: \code{data["model"]}, not \code{data.model}.
Dot access (\code{obj.name}) is for objects such as SDK responses.
\begin{pycode}
>>> d = {"model": "demo-1"}
>>> d["model"]
'demo-1'
>>> d.model
Traceback (most recent call last):
  ...
AttributeError: 'dict' object has no attribute 'model'
\end{pycode}
\end{typecheck}

\begin{boss}[: reading nested data]
Chains of subscripts are read left to right; each one peels off one layer. Always
know the type of the layer you are in: a list needs an \emph{int}, a dict needs a \emph{key}.
\begin{pycode}
>>> resp = {"choices": [{"message": {"text": "ok"}}]}
>>> resp["choices"]
[{'message': {'text': 'ok'}}]
>>> resp["choices"][0]
{'message': {'text': 'ok'}}
>>> resp["choices"][0]["message"]
{'text': 'ok'}
>>> resp["choices"][0]["message"]["text"]
'ok'
\end{pycode}
When one step fails, evaluate the chain one subscript at a time. The error shows
which layer was not what you thought.
\end{boss}

% ------------------------------------------------------------
\section{Review}

You now know the shape of a Python program: names, values, lists, dictionaries,
loops, imports, and indentation for blocks. You know two ways to run code and three
places code comes from. Next chapter: what a variable really is.

\begin{practice}
\begin{enumerate}[itemsep=1pt]
\item Why does \code{models[1]} on a three-item list give the \emph{second} item?
\item What is the difference between what the interpreter echoes and what a file prints?
\item Name the three sources of Python code and give one example of each.
\item Which error do you get from a list offset that is too large? From a dict key that does not exist?
\item In \code{r["a"][0]["b"]}, what type must \code{r}, \code{r["a"]} and \code{r["a"][0]} have?
\item What marks the end of a \code{try} block in Python?
\end{enumerate}
\end{practice}

\begin{achievement}{First Run}
\begin{itemize}[itemsep=1pt]
\item Read a short Python program: loop, list, dictionary, \code{try}/\code{except}.
\item Check your Python version and start the interpreter.
\item Run code interactively and from a \code{.py} file, and explain the echo difference.
\item Tell built-ins, the standard library and third-party packages apart.
\item Navigate nested list/dict data with chained subscripts.
\end{itemize}
\end{achievement}
